% 修士論文中間試問会 Extended Abstract
% A4 2段組・上30/下27/左右18mm・段間7mm・本文10.5pt・表題12pt
\documentclass[a4paper]{ltjsarticle}
\usepackage[top=30mm,bottom=27mm,left=18mm,right=18mm,columnsep=7mm]{geometry}
\usepackage{luatexja}
% 英文はTimes系（要件）
\usepackage{fontspec}
\setmainfont{TeX Gyre Termes}
\usepackage{booktabs}
\usepackage{array}
\usepackage{tikz}
\usetikzlibrary{arrows.meta,positioning}
\usepackage{titlesec}
\usepackage{caption}
\captionsetup{font=small,labelsep=quad}

% 文字サイズ（厳守: 本文10.5pt）
\makeatletter
\renewcommand{\normalsize}{\@setfontsize\normalsize{10.5}{15.5}}
\makeatother
% 見出し10.5pt太字
\titleformat{\section}{\normalsize\bfseries}{\thesection.}{0.6em}{}
\titleformat{\subsection}{\normalsize\bfseries}{\thesubsection}{0.6em}{}
\titlespacing{\section}{0pt}{0.9em}{0.3em}
\titlespacing{\subsection}{0pt}{0.7em}{0.2em}
\pagestyle{empty}

\begin{document}
\normalsize
\twocolumn[{%
\begin{center}
{\fontsize{12}{18}\selectfont\bfseries
対面議論支援のための実時間話者同定とAI介入の設計\par}
\vspace{2pt}
{\fontsize{12}{16}\selectfont
Real-time Speaker Identification and AI Intervention Design\\
for Face-to-face Discussion Support\par}
\vspace{6pt}
{\fontsize{10.5}{14}\selectfont 松井　玲\par}
{\fontsize{10.5}{14}\selectfont Rei Matsui\par}
\end{center}
\vspace{4pt}
\begin{center}
\begin{minipage}{\dimexpr\textwidth-14mm\relax}
\fontsize{10.5}{15}\selectfont
対面の会議や授業では，参加の偏りや論点の散逸が生じやすい．オンラインのテキスト議論ではAIの介入の効果が報告されている一方，重要な議論の多くが行われる対面では，こうした支援はまだ少ない．対面では発話が音声のまま消えていくため，誰が何を言ったかを実時間で把握することが支援の前提となる．しかしこの把握は誤りを避けられず，誤ったまま介入すれば参加者の信頼を損なう．そこで本研究では，相槌ではなく議論の実質的な発話を間違えないことを話者同定の要件と定め，複数の証拠から発話ごとに話者の確信度を求める選択的話者同定を設計・実装した．日本語日常会話コーパスでの評価では，議事録の文字の86.8\%で話者を正しく同定し，名指しの介入を検証できる水準を確認した．修士論文では，この基盤の上で対面の議論においてAIはいつ・誰に・何を介入すべきかを，参加者を伴う実験により検証する．
\end{minipage}
\end{center}
\vspace{10pt}
}]

\let\thefootnote\relax
\footnotetext{\fontsize{10.5}{14}\selectfont 社会情報ネットワーク講座・合意情報学分野 (指導教員: 伊藤孝行)}

\section{はじめに}
会議や授業をはじめとする対面の議論は，参加の偏り・論点の散逸・結論の曖昧さといった構造的な問題を抱えている．これに対し，大規模言語モデル（LLM）を用いて議論に介入するエージェントの研究が活発化しており，オンラインのテキスト議論では，投稿の要約・論点の提示・少数派視点の代弁などの介入が議論の質を高めることが報告されている\cite{ito,gu,fulay,tessler}．

本研究が最終的に目指すのは，こうした支援を対面の議論で実現することである．組織の意思決定，授業，地域の合意形成など，重要な議論の多くはいまも対面で行われている．しかも対面の場では，上下関係や同調圧力が同じ空間で直接働くぶん，参加の偏りを生む力がオンラインより強い．

対面でAIにできる介入は一つではない．第一に，発言の少ない参加者への振りである．発言の機会が参加者に均等に行き渡る集団ほど課題解決の成績が高いことが集団の知能を測る研究で報告されており\cite{woolley}，偏りを減らすことは公平さへの配慮であると同時に，議論の質に関わる．しかしこの働きかけは，気兼ねや上下関係のために参加者どうしでは行いにくく，双方に心理的な負担が生じる．司会を立てても，司会自身が場の利害と関係を持つ限り，この構造は消えない．利害も上下関係も持たない場の外のAIであれば，これを引き受けられる．第二に，論点の整理である．発話が個人に紐づいて保持されていれば，それまでの論点に基づいた介入ができる．例えば，応答のないまま流れた懸念を，誰の発言かという根拠とともに再提示できる．第三に，合意の明確化である．誰の提案が誰の合意を得て決まったのかが名前付きの議事録として残れば，結論の曖昧さや同じ議題の蒸し返しを減らせる．

対面への拡張は，オンラインの手法の単純な移植では成立しない．オンラインでは発言に投稿者のアカウントが自動的に付くため，誰が言ったかはそもそも問題にならない．対面ではこの前提が崩れる．発言は音声のまま消えていき，誰が何を言ったかをシステム自身が実時間で把握しなければ，介入の宛先を決めることも，主張を個人に紐づけることも，参加の偏りを測ることもできない．すなわち，正しい話者同定（誰がいつ何を言ったかを間違えないこと）が対面支援の土台の要件となる．

本研究は二段階で進める．第一段階では，議論の実質的な発話を誰が言ったかを間違えない話者同定を実時間で行う手法を構築し評価する．第二段階では，対面の議論においてAIはいつ・誰に・何を介入すべきかという問いを，構築した基盤を用いた実験により検証する．現在までに完了しているのは第一段階であり，本稿ではその設計と結果を報告し，第二段階の実験計画を示す．介入の内容生成や発言内容の構造化を含む支援システム全体の構築は，この土台の上で今後行う．

\section{関連研究}
\subsection{対面の議論とファシリテーションの効果}
会議中の相互作用の質は，チームの生産性や組織の成果と関連する\cite{kauffeld}．また，発言機会が参加者に均等に行き渡る集団ほど課題解決の成績が高く\cite{woolley}，気兼ねなく発言できるチームほど学習が進むことが知られている\cite{edmondson}．発言の少ない人への促しが参加の偏りを減らすことも，ロボットによる司会の実験で確認されている\cite{asadi}．すなわち，参加の偏りを正す働きかけの効果には裏付けがあり，問題は，その働きかけを日常の会議で誰がどう担うかにある．

\subsection{オンライン議論支援エージェント}
LLM以前から，大規模なオンライン議論をエージェントが進行するD-Agree\cite{ito}が知られ，LLM以降は，議論のファシリテーションを行うPTFA\cite{gu}，熟議に不在の視点をペルソナとして代弁するThe Empty Chair\cite{fulay}，対立する集団の共通基盤の発見を仲介する研究\cite{tessler}など，介入の内容と効果の研究が進んでいる．リモート会議へ対象を広げたCLARA\cite{gunasekaran}も，発言と発言者の対応はプラットフォームから与えられる．さらにテキストのチャット議論では，メッセージの多くを個人宛てに行うLLMの進行が発言の最も少ない参加者の発話を底上げする一方，名指しへの受け止めは分かれることが統制実験で報告されている\cite{alsobay}．すなわちこれらの研究では，話者の同定という問題自体が存在しない．

\subsection{対面の会議支援}
対面を対象とするシステムは，同定の困難を三つの方法で回避してきた．第一に同定を諦める方法である．EchoMind\cite{chen}は対面会議の音声を実時間で認識しLLMで論点地図を構築するが，発話を個人に結び付けず，介入も画面提示に留まる．第二に機材で回避する方法である．Meeting Mediator\cite{kim}は全参加者に装着させたセンサバッジで発話量を測って可視化し，参加の偏りが減ることを示した．商用サービスのHylableも専用のマイクアレイ装置で発話量を可視化する．いずれも文字レベルの話者同定や音声での介入は行わない．第三に人手で回避する方法である．ロボットによる議論の司会では，発話量の均衡化や促しなどの介入の効果が繰り返し確認されているが，その多くは人間の遠隔操作によるもので，話者の判別も人手や個人マイクに頼る構成が主流である\cite{weisswange}．AsadiとFischer\cite{asadi}は，発話時間の追跡により最も寡黙な学生へ発話を促すロボットで参加の偏りが減ることを示したが，判別には人手の監視を併用し，促しは名指しではない．また直近では，マイクアレイ・顔認識カメラ・声と顔の事前登録を組み合わせ，参加者を名前で選んで応答する自律ロボットが報告されているが，音声による同定の正解は群会話の発話の18.4\%に留まり，運用は顔認識に依存した\cite{abbo}．機材と登録を併用しても，音声からの同定自体は解決していない．

まとめると，こうした回避によらず話者を判別し，名指しの音声介入を行う自律システムは，調査した範囲では見当たらない．

\subsection{エージェントの誤りと信頼}
対面で名指しを行う以上，呼び間違いの影響は避けて通れない．人とロボットの相互作用の分野では，ロボットの誤りが利用者の信頼を損なうこと，誤りの種類によって信頼の毀損と回復の仕方が異なることが知られている\cite{salem}．また機械学習の分野では，確信の持てない入力への出力を差し控える棄権（abstention・selective prediction）が，LLMの誤答抑制の文脈で活発に研究されている\cite{wen}．ただし，その対象は主に単発の質問応答であり，実時間の多人数対話の振る舞いとして棄権を設計・検証した例は乏しい．

\subsection{話者分離と話者付き書き起こし}
音声処理分野では話者分離（diarization）と話者付き書き起こし（speaker-attributed ASR）の研究が進む．ストリーミング話者分離の代表的手法であるLS-EEND\cite{liang}は会議コーパスAMIでDER 20.8\%を達成するが，これは対象ドメインでの追加学習を経た値である．同種の実時間分離としてStreaming Sortformer\cite{medennikov}も提案されているが，扱える話者は最大4人で，出力は到着順の匿名ラベルである．話者付き書き起こしのDNCASR\cite{zheng}はオフライン処理を前提とする．いずれも出力は匿名の話者ラベルであり，介入に必要な実名の同定までは扱わない．

\section{対面AI介入の要件と本研究の問い}
\subsection{対面がオンラインと異なる3条件}
対面の議論には，オンラインに存在しない3つの条件がある．（1）介入が進行中の議論に割り込む．音声の介入は読み飛ばせず，画面への提示も参加者の注意を議論から引き離すため，タイミングを誤ると議論そのものを妨げる．（2）発言が消えていく．画面上のログがなく，システムが誰の発言かを保持しなければ介入の根拠を示せない．（3）名指しに社会的コストがある．宛先を持つ介入は効果が大きい一方，呼び間違いは参加者の信頼を大きく損なう．

本研究はこの環境を，マイク1本という最小の機材構成で扱う．参加者ごとの個人マイクや専用アレイの設置・装着は導入の障壁となるのに対し，手元のPCを置くだけで始められる構成は日常の会議に持ち込みやすい．また，最も制約の強い構成で同定が成立すれば，個人マイクなどを利用できる環境では同等以上の性能が期待できる．なお本設計は機材の追加を排除せず，個人マイクが利用できる場ではそれを証拠源の一つとして手法に加えられる．

\subsection{話者同定に求める要件}
条件（3）から，話者同定に求められるのは全発話の完全な同定ではないことが導かれる．介入が名指しの根拠とするのは議論の実質的な発話であり，長い主張の同定を誤れば誤った名指しに直結する．一方，相槌のような短い反応も，無理に同定して誤れば，同意や反応を別人のものとして記録に残してしまう．そこで本研究は話者同定の要件を「議論の実質的な発話を，誰が言ったかを間違えないこと」と定義し，これを3つの設計判断として具体化した．（a）評価は発話数ではなく文字数で重み付けし，相槌は対象外とする．（b）確信が持てない発話には名前を出さず「未確定」とする．（c）介入の名指しは実質的な発話への応答に限る．

\subsection{介入について検証する問い}
この基盤の上で検証する問いは「対面の議論において，AIはいつ・誰に・何を介入すべきか」である．この問いは3軸に分解できる．いつ（どのような沈黙・停滞で介入するか），誰に（全体に言うか，特定の参加者を名指しするか），何を（要約か，論点の提示か，発言の少ない人への振りか）である．

この3軸の分解は，多人数対話エージェントのサーベイ\cite{sapkota}が社会的なエージェントの基本能力を「いつ話すか・誰に向けるか・何を話すか」として整理していることと対応する．ただし同サーベイが扱う研究の大半はテキストないし統制された対話環境を前提としており，対面の実環境で「誰に」を実現する手段は未解決のまま残っている．

本研究の実験（第8節）はまず「誰に」の軸を切り出す．理由は2つある．第一に，この軸は話者同定の価値が最も直接に試される軸であり，同定の評価と介入の検証を一つの実験で行える．第二に，「いつ」と「何を」は議論の内容や課題に依存するため統制が難しいのに対し，「誰に」は介入文の宛先だけを操作すればよく，内容を固定した比較が成立する．残る2軸は，実験で得られる観察データ（どの介入が応答され，どれが無視されたか）から次の問いを立てる．

\section{提案手法: 選択的話者同定}
\subsection{手法の考え方}
提案する選択的話者同定は，発話ごとに複数の独立した証拠から話者の確信度を求め，確信を持てた発話に名前を付ける手法である．確信を持てなかった発話は未確定として区別する．手法が入力として要求するのは，(i) 発話の文字と仮の話者ラベル，(ii) 声の区別（どの発話が同じ声か），(iii) 声の照合（過去の誰の声に似ているか）という3種類の証拠であり，これらを与える認識器は交換できる．本稿の実装では (i) に音声認識のSoniox，(ii) にストリーミング話者分離のpyannote，(iii) に声紋モデルのReDimNet\cite{yakovlev}を用いる（認識器を交換した場合の影響は第5.4節で確認する）．なお本稿の話者同定は，登録済みの話者集合から声の主を選ぶ古典的な話者同定とは異なり，話者の集合を事前に仮定せず，会議の進行の中で話者を獲得しながら発話単位で名前を対応付ける．

この設計は，AIが介入するという用途から導かれている．書き起こしの誤りは後から直せるが，介入はその時点の判定で宛先を決め，一度口にした名指しは取り消せない．したがって名指しは，確信を持てた同定に基づいて行う必要がある．介入が参照する議論の実質的な発話は長く，証拠が揃いやすいため，この要求と両立する．逆に相槌のような短い発話は，誤って同定すると同意や反応が別人の記録として残るため，無理に判定せず未確定とする．図\ref{fig:arch}に構成を示す．3種類の証拠はそれぞれ異なる誤り方をする．文字のラベルは話者交替に弱く，声の区別は同一人物を複数のクラスタに割り，声の照合は短い発話で不安定である．手法はこれらを突き合わせ，複数の証拠が一致する発話に確信を割り当てる．

\begin{figure}[tb]
\centering
\begin{tikzpicture}[
  box/.style={draw,rounded corners=2pt,align=center,font=\small,
              minimum height=8mm,inner sep=3pt},
  nav/.style={box,fill=black!75,text=white},
  arr/.style={-{Stealth[length=2mm]},semithick}]
\node[box,minimum width=78mm] (mic) {音声入力（マイク1本）};
\node[box,below=5mm of mic,minimum width=25mm]
  (dia) {話者分離};
\node[box,left=1.5mm of dia,minimum width=25mm]
  (stt) {音声認識};
\node[box,right=1.5mm of dia,minimum width=25mm]
  (vp) {声紋照合};
\node[nav,below=5mm of dia,minimum width=78mm]
  (med) {選択的話者同定（本研究で実装）\\[-2pt]{\footnotesize 話者枠の割当て・遡及訂正・判定の保留}};
\node[box,below=5mm of med,minimum width=78mm]
  (out) {発話ごとの話者（名前または未確定）};
\node[below=4mm of out,xshift=-19mm,font=\small] (minutes) {議事録};
\node[below=4mm of out,xshift=19mm,font=\small] (interv) {介入層};
\draw[arr] (mic.south -| stt) -- (stt.north);
\draw[arr] (mic.south -| dia) -- (dia.north);
\draw[arr] (mic.south -| vp) -- (vp.north);
\draw[arr] (stt.south) -- (med.north -| stt);
\draw[arr] (dia.south) -- (med.north -| dia);
\draw[arr] (vp.south) -- (med.north -| vp);
\draw[arr] (med) -- (out);
\draw[arr] (out.south -| minutes) -- (minutes.north);
\draw[arr] (out.south -| interv) -- (interv.north);
\end{tikzpicture}
\caption{提案手法の構成}
\label{fig:arch}
\end{figure}

\subsection{手法の構成}
手法は3つの処理からなる．

\textbf{話者枠の割当て．}参加人数ぶんの話者枠を用意し，枠ごとにその人の声の特徴を蓄積する．証拠だけでは決めきれなかった発話は，枠に貯めた声と照合し，コサイン類似度が受理しきい値以上なら割り当てる．話者分離が同一人物を複数クラスタに割った場合も，ここで一つにまとめる．

\textbf{遡及訂正．}序盤は参照できる声紋が少なく判定が安定しない．そこで全発話の声紋を保持し，2分・5分・以降2分ごとに，その時点の参照で過去の判定をやり直して記録に反映する．介入は過去の発話の記録を参照して内容を組み立てるため，この訂正により，会議が進むほど介入は正しい文脈に基づけるようになる．議事録の質も同じ訂正で保たれる（発話直後の判定の精度は第5.5節で別途測る）．

\textbf{判定の保留．}相槌のような短すぎる発話は判定の対象としない．候補どうしの類似度の差が僅かなときや，既知の誰とも似ない声のときは，未確定として出力する．

話者枠の呼称は，はじめは参加者A・Bのような匿名のラベルであり，実名が得られない参加者は匿名のまま扱う．実名との対応づけの一つとして，自己紹介の発話からその声と実名を結びつける名乗りの検出を補助的に実装している（予備的な確認では，記者会見動画の自己紹介5件をすべて検出し，80回の実行で誤検出はない）．

声紋照合の受理しきい値・僅差とみなす幅・遡及訂正のタイミングの3つは，第5.1節の開発セット9会話で調整した（受理はコサイン類似度0.42以上，僅差は1秒未満の発話で候補間の差が0.03未満）．テストセットとコーパス外の音声には，この値を一切変更せず持ち込んでいる．

\subsection{実時間性と表示の安定}
システムはノートPC1台で実時間に動作し，議事録は発話とほぼ同時に画面へ表示される．遡及訂正による話者名の変更は個々の発話の貼り直しとして現れ，一度表示した人物の呼称は会議を通じて固定されるため，利用者が誰が誰かを見失わない．介入の音声再生区間も議事録と同じ時計で記録し，介入が人の発話と重なっていないかを事後に検証できるようにした（第6.4節）．

\subsection{介入層への受け渡し}
発話ごとに，未確定かどうか・確信度・判定の経路を介入層へ渡し，未確定の相手は名指ししない．参加度の集計は，同定された名前ごとの回数ではなく話者分離が測った発話時間で行い，未確定の増減が参加度の見積りを歪めないようにした．

\section{評価}
\subsection{データと指標}
正解付きの評価には千葉大学の3人日常会話コーパス\cite{chiba}から13会話を用いた．うち9会話をしきい値調整に用いる開発セット，4会話を調整に使わないテストセットとし，評価はテストセットを基準に読む（開発セットはしきい値を合わせたぶん高めに出る）．評価は，コーパスの音源を実時間で再生してシステムに入力する再生ランとして行った．実際の会議室で収音した音声による検証は今後行う（第8節）．

指標は次のとおり定義した．\textbf{正解率}＝全文字のうち話者を正しく同定できた文字の割合，\textbf{誤同定率}＝別人と同定した文字の割合であり，残りが\textbf{未確定}（名前を出さなかった分）である．発話数ではなく文字数で数えるのは，長い発言の誤りほど重く数えるためである．相槌は分母から除いた．システムの発話と正解の対応付けは，時間の重なりではなく文章の一致（正規化した文字列の類似）で行う．発話の区切りが正解と一対一に対応せず，また発話時間の24.8\%が重なっているため，時間による対応付けでは採点できる範囲が6割弱に留まるのに対し，文章の一致では84\%を採点できる．

未確定を誤同定と区別して数えるのは，未確定を本設計が正しい挙動として扱っているためである．両者を混ぜた指標では，誤った名指しを避けて黙るシステムと黙らず誤るシステムの区別がつかず，本研究の設計判断そのものが評価できない．また，文章の一致による対応付けは音声認識の誤りに対して頑健である．認識が大きく崩れた発話は類似度がしきい値を下回り採点対象外へ落ちるため，誤った正解と突き合わせて誤った点を付けることがない．実際に対応が付かないのは平均0.9秒の相槌が大半である．

\subsection{同定の精度}
表\ref{tab:main}に主結果を示す．人数を与えた通常の使い方で，テストセットの正解は86.8\%，誤同定は5.8\%である．名前を出した発話に限れば93.7\%が正しく，未確定は7.4\%であった．人数を与えない場合も，低下分は主に未確定の増加であり誤同定はほとんど増えない．すなわち条件が悪化したとき，増えるのは誤りではなく未確定である．これは第3.2節の要件（b）が設計どおり機能していることを示す．

\begin{table}[tb]
\centering
\caption{同定の精度（文字ベース，正解率／誤同定率）}
\label{tab:main}
\small
\begin{tabular}{lcc}
\toprule
条件 & 開発9会話 & テスト4会話 \\
\midrule
人数を与える & 91.5\% ／ 4.3\% & 86.8\% ／ 5.8\% \\
人数を与えない & 85.5\% ／ 4.5\% & 79.5\% ／ 4.7\% \\
\bottomrule
\end{tabular}
\end{table}

\subsection{名指しの被覆と正しさ}
同定の内訳を，介入の観点から「名前を出す割合（被覆）」と「名前を出したときの正しさ」に組み替えたものを表\ref{tab:risk}に示す．条件が悪化しても動くのは被覆であり，名前を出したときの正しさは93.7--95.5\%の狭い帯に留まる．本設計の下では，環境の悪化は名指しの質ではなく名指しの量に現れる．

\begin{table}[tb]
\centering
\caption{名指しの被覆と正しさ（文字ベース）}
\label{tab:risk}
\small
\begin{tabular}{lcc}
\toprule
条件 & 名前を出す被覆 & 名指し中の正解 \\
\midrule
人数あり・開発 & 95.8\% & 95.5\% \\
人数あり・テスト & 92.6\% & 93.7\% \\
人数なし・開発 & 90.0\% & 95.0\% \\
人数なし・テスト & 84.2\% & 94.4\% \\
\bottomrule
\end{tabular}
\end{table}

\subsection{寄与の分解}
性能の内訳を，同じ開発セットで構成を段階的に加えながら測った（表\ref{tab:decomp}）．音声認識の話者ラベルをそのまま使うと47.9\%，話者分離と声紋を加えた既製の認識器の総合でも67.0\%に留まるのに対し，提案手法を加えると91.5\%まで向上する．また，声紋モデルを別の公開モデルに交換しても全体の正解率は91.5\%から91.9\%とほとんど変わらず，話者分離器の交換でも変わらなかった．ただし，試した代替は元の認識器と同水準の性能のものであり，より高性能な認識器に置き換えた場合の検証は今後の課題である．

\begin{table}[tb]
\centering
\caption{寄与の分解（開発セット9会話・文字ベース正解率）}
\label{tab:decomp}
\small
\begin{tabular}{lc}
\toprule
構成 & 正解率 \\
\midrule
音声認識の話者ラベルのみ & 47.9\% \\
＋話者分離＋声紋（既製の認識器の総合） & 67.0\% \\
＋選択的話者同定（本研究で実装） & 91.5\% \\
\bottomrule
\end{tabular}
\end{table}

\subsection{介入時点の精度}
議事録の最終値は遡及訂正を経た値だが，発話の直後に行う介入の宛先決定は，訂正を待たずその時点の判定を使う．そこで遡及訂正を行う前の判定を，同じ実装・同じ採点で測った（表\ref{tab:now}）．発話直後の実時間の正解率は開発84.7\%・テスト86.0\%である．遡及訂正は開発ではこれを91.5\%まで引き上げ，テストでは86.8\%との差は小さかった．

\begin{table}[tb]
\centering
\caption{介入時点の精度（文字ベース正解率）}
\label{tab:now}
\small
\begin{tabular}{lcc}
\toprule
時点 & 開発9会話 & テスト4会話 \\
\midrule
遡及訂正の前（発話直後の判定） & 84.7\% & 86.0\% \\
遡及訂正の後（議事録の最終値） & 91.5\% & 86.8\% \\
\bottomrule
\end{tabular}
\end{table}
経過時間帯別の値を表\ref{tab:band}に示す．開始0--1分は正解が58--61\%まで落ちるが，誤同定は3--5\%に留まり，大半は未確定として現れる．精度は時間とともに上がり，5分以降は92\%前後に達する．

\begin{table}[tb]
\centering
\caption{経過時間帯別の発話直後の精度（文字ベース，正解率／誤同定率）}
\label{tab:band}
\small
\begin{tabular}{lcc}
\toprule
経過時間 & 開発9会話 & テスト4会話 \\
\midrule
0--1分 & 61.0\% ／ 5.3\% & 58.1\% ／ 3.4\% \\
1--2分 & 74.0\% ／ 16.2\% & 81.4\% ／ 3.9\% \\
2--5分 & 85.6\% ／ 11.4\% & 87.5\% ／ 7.7\% \\
5分以降 & 92.5\% ／ 4.9\% & 91.9\% ／ 5.7\% \\
\bottomrule
\end{tabular}
\end{table}

\subsection{標準指標と英語音声での確認}
日本語13会話の標準指標はDER 39.5\%（内訳: 検出漏れ19.5，誤検出2.6，取り違え17.3．未確定を除いた取り違えは7.5\%），cpCER 63.3\%（うち音声認識単体の文字誤りが28.8\%）であった．誤差の主因は同時発話の検出漏れであり，取り違えではない．

英語会議コーパスAMI\cite{ami}の1会議でも同一システム（追加学習なしのゼロショット）を測定し，DER 45.1\%（話者を検出できた区間に限れば32.7\%），未確定を除く取り違え6.7\%，cpWER 53.8\%であった．取り違えの水準が日本語（7.5\%）と英語（6.7\%）でほぼ同じであることから，この構成の同定は言語に強く依存しないことが示唆される．ただし英語は1会議での確認である．なおAMIで追加学習したストリーミング手法LS-EEND\cite{liang}のDER 20.8\%とは差があるが，主因は同時発話の検出漏れであり，取り違えの水準は同等である．条件を揃えた比較は今後の課題とする．

\subsection{コーパス外の音声での確認}
しきい値の過適合を確認するため，調整に一切使っていない公開動画4本で挙動を目視で確認した（表\ref{tab:wild}）．コーパス外の未知の話者でも，討論や雑談の音声では話者の取り違えは見られず，4人の速い割り込みでも話者枠は最後まで安定した．例外は記者会見で，遠距離マイクと残響により別人の声が似て，複数の記者が同じ話者枠に統合された．自己紹介からの登録（名乗りの検出）で部分的には補えたが，この環境（多数の一言話者×遠距離マイク）は現在の構成の適用範囲外と整理している．なおこれらは正解付けをした測定ではなく目視の抜き取り確認であり，実際の会議室での検証は今後行う（第8節）．

\begin{table}[tb]
\centering
\caption{コーパス外の公開動画4本での確認（目視）}
\label{tab:wild}
\small
\begin{tabular}{lp{30mm}p{26mm}}
\toprule
素材 & 条件 & 結果 \\
\midrule
討論番組3人 & スタジオ収録・議論の応酬 & 取り違えなし \\
討論番組4人 & 速い割り込みが多い & 4枠とも安定 \\
収録雑談3人 & クリーンな収録 & 崩れなし \\
記者会見 & 遠距離マイク・多数の一言話者 & 別人が同一枠に統合 \\
\bottomrule
\end{tabular}
\end{table}

\section{分析}
\subsection{機構別の寄与}
提案手法の各機構を個別に外した対照を表\ref{tab:ablation}にまとめる．遡及訂正を外すと正解率は91.5\%から84.7\%へ下がり，誤同定は4.3\%から8.4\%へ増える．人数情報（話者枠の割当ての入口）を外すと67.5\%まで落ち，これは既製の認識器の総合67.0\%とほぼ一致する．すなわち提案手法の寄与の大半は話者枠の割当てとその周辺にある．なお表\ref{tab:main}の「人数を与えない」は話者枠の機構を保ったまま人数の指定だけを外した運用条件であり，機構自体を外した本対照とは異なる．一方，判定の保留（短い発話の棄権）は正解率をほとんど変えないまま，発話数ベースの誤同定を12.6\%から9.2\%へ削る．保留は正解率を高める機構ではなく，誤りを未確定に置き換える機構であることが分かる．

\begin{table}[tb]
\centering
\caption{機構別の対照（開発セット・文字ベース）}
\label{tab:ablation}
\small
\begin{tabular}{lcc}
\toprule
条件 & 正解率 & 誤同定 \\
\midrule
提案手法（全機構） & 91.5\% & 4.3\% \\
遡及訂正を外す & 84.7\% & 8.4\% \\
話者枠の上限を外す & 85.8\% & 4.5\% \\
人数情報をすべて外す & 67.5\% & 13.9\% \\
\bottomrule
\end{tabular}
\end{table}

\subsection{残る誤同定の内訳}
残った誤同定を，発話ごとの判定経路別に分解した（正解付き13会話の再集計，表\ref{tab:kind}）．誤同定の9割以上は，声紋照合が高い確信を持てないまま他の手がかり（音声認識のラベル継続・クラスタの混在・声紋の蓄積不足）で決めた発話から生じており，声紋が高信頼で一致した発話からの誤りは2\%程度に過ぎない．すなわち残る誤りの本質は照合器の精度不足ではなく，自信のない発話をどう扱うかという問題である．この分解が，しきい値の引き上げ（照合の強化）ではなく棄権の拡充が正しい改善であることの根拠となった．また発話長で見ると，誤同定は1秒未満の短い発話の帯に集中しており，棄権規則の導入によりこの帯の誤同定は40\%から14.6\%へ低下した．

\begin{table}[tb]
\centering
\caption{誤同定の判定経路別内訳（13会話）}
\label{tab:kind}
\small
\begin{tabular}{lc}
\toprule
判定経路 & 誤同定に占める割合 \\
\midrule
ラベル混在（音声認識の話者混載） & 43.7\% \\
ラベル継続（前話者の引き継ぎ） & 29.2\% \\
声紋の蓄積不足 & 19.5\% \\
補正・高信頼一致・その他 & 7.6\% \\
\bottomrule
\end{tabular}
\end{table}

\subsection{誤り注入の感度分析}
この分析の目的は2つある．第一に，現在の同定精度のまま名指しの介入を行うと，呼び間違いがどの程度起きるかを見積もること．第二に，同定が今より悪化したとき，介入がどれだけ壊れるかを知ることである．そのために，同定結果の名前を確率$\varepsilon$で別の参加者に置き換えるシミュレーションを行った（正解付き13会話，200試行平均）．表の$\varepsilon=0$の行が現行システムの実測に相当し，下の行ほど同定を悪化させた場合を表す．介入が同定を使う場面としては，直近の発話者を名指しして応答する介入（名指し）と，発話量が最少の参加者に話を振る介入（声かけ）を模擬した．実際の誤りは音響的に似た特定の2人に偏るため，一様な置き換えは楽観側の近似である．

表\ref{tab:inject}に開発セットの結果を示す．第一に，発話を選ばず名指しすると現行（$\varepsilon=0$）でも21\%が誤るのに対し，20文字以上の実質発話への応答に限れば3.1\%（テストセット6.6\%）まで下がる．誤同定が短い発話に集中していることの帰結であり，第3.2節の要件（c）の定量的な根拠である．第二に，注入した誤りはほぼそのまま名指しの誤りに乗る（誤りを5\%注入すると名指しの誤りも約5\%増える）．したがって，名指しの誤りを$x$\%以下にするという要求から，同定精度の要件を逆算できる．第三に，発話量最少者の特定は集計量のため悪化は緩やかだが，発話量が拮抗する場面では現行でも誤りうる．声かけの介入文は最少の人と断定しない形が安全である．なお声かけの列は判定時点の数が少なく（開発27点・テスト12点），粗い見積りである．

\begin{table}[tb]
\centering
\caption{誤り注入の感度分析（開発セット）}
\label{tab:inject}
\small
\begin{tabular}{cccc}
\toprule
注入$\varepsilon$ & \begin{tabular}{c}名指し誤り\\（全発話）\end{tabular} &
\begin{tabular}{c}名指し誤り\\（20字以上）\end{tabular} &
\begin{tabular}{c}声かけ相手\\の誤り\end{tabular} \\
\midrule
0\%（現行） & 21.1\% & 3.1\% & 18.5\% \\
+5\% & 24.5\% & 8.0\% & 24.0\% \\
+15\% & 31.2\% & 17.0\% & 27.3\% \\
+30\% & 41.7\% & 31.9\% & 35.4\% \\
\bottomrule
\end{tabular}
\end{table}

テストセットでも傾向は同じで（表\ref{tab:inject2}），20字以上の名指し誤りは$\varepsilon=0$で6.6\%，注入分がほぼそのまま乗る．

\begin{table}[tb]
\centering
\caption{誤り注入の感度分析（テストセット）}
\label{tab:inject2}
\small
\begin{tabular}{cccc}
\toprule
注入$\varepsilon$ & \begin{tabular}{c}名指し誤り\\（全発話）\end{tabular} &
\begin{tabular}{c}名指し誤り\\（20字以上）\end{tabular} &
\begin{tabular}{c}声かけ相手\\の誤り\end{tabular} \\
\midrule
0\%（現行） & 25.6\% & 6.6\% & 0.0\% \\
+5\% & 28.6\% & 10.9\% & 11.7\% \\
+15\% & 34.8\% & 20.2\% & 16.4\% \\
+30\% & 44.1\% & 33.8\% & 26.6\% \\
\bottomrule
\end{tabular}
\end{table}
また，名指しを実質発話への応答に限る条件では，自然発生の呼び間違いは1セッションあたり0.3--0.8回と稀であることも見積もられた．誤りのコスト自体を実験で測るには，制御された提示（視聴実験）が必要という根拠になる．

\subsection{介入側の動作確認}
介入機構の動作も点検した（表\ref{tab:audit}）．12回の実行・発火73件において，人の発話への割り込み（要求された沈黙の間の違反）は0件，同じ言い回しの繰り返しは0組，介入の頻度は10分あたり中央値2.9回（間隔の中央値91秒）であり，介入の実行機構としては健全に動作している．介入の内容が議論の役に立つかは未検証であり，これが第二段階で検証する問いである．

\begin{table}[tb]
\centering
\caption{介入機構の点検（12実行・発火73件）}
\label{tab:audit}
\small
\begin{tabular}{ll}
\toprule
観点 & 結果 \\
\midrule
人の発話への割り込み & 違反 0 / 60 件 \\
同じ言い回しの繰り返し & 0 / 221 組 \\
介入の頻度 & 10分あたり中央値 2.9 回 \\
介入の間隔 & 中央値 91 秒 \\
\bottomrule
\end{tabular}
\end{table}

\section{考察}
本研究の設計の中心は，介入に使う同定を，複数の証拠が一致し確信を持てた発話に限ることである．評価の結果はこの設計が機能していることを示している．条件の悪化は名指しの正しさ（93--96\%）ではなく被覆に現れ（表\ref{tab:risk}），棄権は正解率を犠牲にせず誤同定だけを削り（第6.1節），序盤の弱さも呼び間違いではなく未確定として現れる（第5.5節）．誤りのコストが非対称な社会的場面で動くエージェントに対して，この設計方針は再利用できる可能性がある．

感度分析（第6.3節）では，同定の誤りと名指しの誤りの関係が定量的に分かった．注入した誤りがほぼ1:1で名指しの誤りに乗るため，名指しの誤りを$x$\%以下にしたいという介入側の要求から，同定側に必要な精度を逆算できる．この逆算は，今後の要件定義に使える．

提案手法は既製の音声認識・話者分離・声紋照合を証拠源として用いるが，寄与の分解（表\ref{tab:decomp}）が示すとおり，認識器の総合は67.0\%に留まり，全体の91.5\%との差は認識器では説明できない．認識器を交換しても性能がほとんど動かないことも併せると，本研究の実体は認識器ではなく，異なる誤り方をする証拠を統合し，確信が持てないときは保留するという手法の側にある．また，認識器の統合が進み精度が向上しても，この層の必要は残ると考えている．書き起こしを目的とする製品では，迷った場合も最尤の話者を付けて出力する方が議事録として有用であり，確信度を較正して判定を保留する動機がない．介入では逆に，保留が正しい挙動になる．認識器側の精度向上は，感度分析（第6.3節）の$\varepsilon$を下げる改善として本手法にそのまま取り込まれる．

構築したシステムは研究の最終成果物ではなく，介入の問いに答えるための実験装置である．対面議論への介入という問いは，これまで検証の装置が存在しなかったために，オンラインでの代替実験か，人間が裏で操作する疑似実験でしか扱えなかった．誰の発話かを実時間で判別し，宛先を持つ介入を自律的に行える装置が整ったことで，宛先の効果や誤りのコストといった介入設計の変数を，対面の実環境で初めて実験によって検証できるようになる．

全二重（聞きながら話す）の音声対話モデルが登場し，相槌や割り込みの判断を実時間で行う対話が可能になりつつある．これらのAPIが利用可能になれば，本システムの介入の発話生成と割り込み制御は大きく簡略化されるだろう．一方で，本研究が扱う誰が話したかの判別と，誤りうる同定の上での宛先の決め方は，対話モデルの進歩では解決されない独立の問題であり，むしろ介入の表現力が上がるほど宛先の正しさが問われるようになる．本研究の話者同定と設計規則は，その世代のシステムにもそのまま持ち越せると考えている．

最後に，本研究の限界を述べる．第一に，同時発話（発話時間の24.8\%）が誤差の主因として残っており，これは音源分離という別の手段を要する．第二に，遠距離マイク環境（記者会見型）は適用範囲外である．第三に，評価はコーパスの再生と公開動画に依存しており，実際の会議室での検証をこれから行う．第四に，感度分析はシミュレーションであり，名指しの誤りが人間の信頼に与える実際の影響は次節の実験で測る．第五に，視聴実験は第三者視点の評定であり，名指しされる当事者としての受け止めはライブ実験でのみ測れる．

\section{今後の展望と修士論文作成までの予定}
\subsection{検証の設計}
「いつ・誰に・何を」の3軸のうち，まず「誰に」を2つの実験で検証する．

\textbf{視聴実験（10月）．}録音済みの3人議論から2--3分の区間を6箇所切り出し，各区間の末尾にAIの介入音声を合成で重ねる．条件は（A）直前の発話者への正しい名指し・（B）同一内容で宛先なし・（C）別の参加者名による誤った名指しの3水準で，参加者（10--15名）はランダムな組の刺激を視聴し，介入ごとに有用さ・タイミング・「このAIは議論を正しく把握しているか」・受容を7件法で評定する．分析はAとBの対比で宛先の効果を，AとCの対比で誤名指し1回のコストを，Cの後続刺激への持ち越しで信頼毀損の持続を見る．内容を固定して宛先だけを操作するため各効果を分離でき，誤名指しは合成による提示のため参加者を欺かずに扱える（C条件は終了時に開示する）．

\textbf{対面参加者実験（10--11月）．}3人1グループの合意形成課題（30--40分）で，同一内容の介入を名指しあり／宛先なしで無作為に混在させる被験者内比較を行う．仮説はH1「名指し介入は宛先なし介入より応答されやすい」，H2「発言の少ない人への名指しの振りは，宛先なしの振りより，その人のその後の発話量を増やす」である．手順は，同意取得の後，システムが参加者の声と名前を得る練習議論5分，本議論30--40分（介入あり），介入箇所を頭出しして提示するリプレイ評定と質問紙，の順とする．測定は，行動指標として介入後10秒の応答の有無と応答者・振りの対象者の前後5分の発話時間（いずれも実装済みのログから自動取得），主観指標として介入ごとの適切さ評定と自動化への信頼尺度を用いる．名指しは感度分析の設計規則に従い実質発話への応答に限定する．この条件下で自然発生の呼び間違いは1セッション1回未満のため，本実験に欺瞞は含まれない．応答率の差をA 80\%対B 50\%と仮定した検出力分析（$\alpha=.05$，検出力0.8）から各条件約40介入，8グループ24名を目安とする．なお，発言の少ない人への振りが参加の偏りを減らすことはロボット司会の先行研究\cite{asadi}が，個人宛ての介入が発言の少ない参加者の発話を底上げする一方で名指しの受け止めは分かれることはテキスト議論の統制実験\cite{alsobay}が示しており，本実験の新規性は名指しという宛先の形と，その宛先をシステム自身の話者判別に基づいて決める点の検証にある．

\subsection{スケジュール}
表\ref{tab:sched}に修士論文作成までの計画を示す．

\begin{table}[tb]
\centering
\caption{修士論文作成までの予定}
\label{tab:sched}
\small
\begin{tabular}{lp{52mm}}
\toprule
時期 & 内容 \\
\midrule
9月 & 実会議（ゼミ等）の録音と同一手順での同定の測定．視聴実験の刺激作成と倫理手続き \\
10月 & 視聴実験の実施．介入の宛先を話者同定で決める実装 \\
10--11月 & 対面参加者実験（名指しあり／なし） \\
11--12月 & 分析・考察，修士論文の執筆 \\
\bottomrule
\end{tabular}
\end{table}

\begin{thebibliography}{99}
\fontsize{9}{10.6}\selectfont
\bibitem{ito} Ito, T., et al.:
An Agent that Facilitates Crowd Discussion,
\textit{Group Decision and Negotiation}, Vol.~31 (2022).
\bibitem{gu} Gu, W., et al.:
PTFA: An LLM-Based Agent for Discussion Facilitation,
\textit{Proc. PRICAI 2025} (2025).
\bibitem{fulay} Fulay, S. and Roy, D.:
The Empty Chair: Using LLMs to Raise Missing Perspectives in Policy Deliberations,
\textit{Proc. ACM CSCW 2025} (2025).
\bibitem{tessler} Tessler, M.~H., et al.:
AI can help humans find common ground,
\textit{Science}, Vol.~386 (2024).
\bibitem{woolley} Woolley, A.~W., Chabris, C.~F., Pentland, A., Hashmi, N. and Malone, T.~W.:
Evidence for a Collective Intelligence Factor in the Performance of Human Groups,
\textit{Science}, Vol.~330 (2010).
\bibitem{kauffeld} Kauffeld, S. and Lehmann-Willenbrock, N.:
Meetings Matter: Effects of Team Meetings on Team and Organizational Success,
\textit{Small Group Research}, Vol.~43 (2012).
\bibitem{edmondson} Edmondson, A.:
Psychological Safety and Learning Behavior in Work Teams,
\textit{Administrative Science Quarterly}, Vol.~44 (1999).
\bibitem{chen} Chen, W., et al.:
EchoMind: Supporting Real-time Complex Problem Discussions through Human-AI Collaborative Facilitation,
\textit{Proc. ACM Hum.-Comput. Interact. (CSCW)} (2025).
\bibitem{gunasekaran} Gunasekaran, T.~S., et al.:
CLARA: AI-Mediated Facilitation for Enhancing Group Cognition and Cohesion in Remote Collaboration,
\textit{ACM Trans. Comput.-Hum. Interact.} (2025).
\bibitem{alsobay} Alsobay, M., Rothschild, D.~M., Hofman, J.~M. and Goldstein, D.~G.:
Bringing Everyone to the Table: An Experimental Study of LLM-Facilitated Group Decision Making,
\textit{Proc. ACM Hum.-Comput. Interact. (CSCW)} (2026).
\bibitem{kim} Kim, T., et al.:
Meeting Mediator: Enhancing Group Collaboration using Sociometric Feedback,
\textit{Proc. ACM CSCW 2008} (2008).
\bibitem{weisswange} Weisswange, T.~H., Javed, H., Dietrich, M., Jung, M.~F. and Jamali, N.:
Design Implications for Robots that Facilitate Groups---A Scoping Review on Improving Group Interactions through Directed Robot Action,
\textit{ACM Trans. Human-Robot Interaction} (2025).
\bibitem{asadi} Asadi, A. and Fischer, K.:
Not just another quiet student: reducing participation imbalance through robot moderation,
\textit{Frontiers in Education}, Vol.~10 (2025).
\bibitem{abbo} Abbo, G.~A., Pinto-Bernal, M.~J., Catrycke, M. and Belpaeme, T.:
Multi-party open-ended conversation with a social robot,
\textit{Frontiers in Robotics and AI}, Vol.~13 (2026).
\bibitem{salem} Salem, M., Lakatos, G., Amirabdollahian, F. and Dautenhahn, K.:
Would You Trust a (Faulty) Robot? Effects of Error, Task Type and Personality on Human-Robot Cooperation and Trust,
\textit{Proc. ACM/IEEE HRI 2015} (2015).
\bibitem{sapkota} Sapkota, S., Hasan, M.~S., Shah, M. and Karmaker, S.:
Multi-Party Conversational Agents: A Survey,
arXiv:2505.18845 (2025).
\bibitem{wen} Wen, B., et al.:
Know Your Limits: A Survey of Abstention in Large Language Models,
\textit{Trans. ACL} (2025).
\bibitem{yakovlev} Yakovlev, I., et al.:
Reshape Dimensions Network for Speaker Recognition,
\textit{Proc. INTERSPEECH 2024} (2024).
\bibitem{liang} Liang, D. and Li, X.:
LS-EEND: Long-Form Streaming End-to-End Neural Diarization with Online Attractor Extraction,
\textit{IEEE/ACM Trans. Audio, Speech, and Language Processing} (2025).
\bibitem{medennikov} Medennikov, I., Park, T., Wang, W., Huang, H., Dhawan, K., Wang, J., Balam, J. and Ginsburg, B.:
Streaming Sortformer: Speaker Cache-Based Online Speaker Diarization with Arrival-Time Ordering,
\textit{Proc. INTERSPEECH 2025} (2025).
\bibitem{zheng} Zheng, X., Zhang, C. and Woodland, P.:
DNCASR: End-to-End Training for Speaker-Attributed ASR,
\textit{Proc. ACL 2025} (2025).
\bibitem{chiba} 伝康晴・榎本美香:
『千葉大学3人会話コーパス』使用説明書 Release 1,
国立情報学研究所 音声資源コンソーシアム (2014).
\bibitem{ami} Carletta, J.:
Unleashing the killer corpus: experiences in creating the multi-everything AMI Meeting Corpus,
\textit{Language Resources and Evaluation}, Vol.~41 (2007).
\end{thebibliography}

\end{document}
